\subsubsection{k3s}

Da der Production Cluster im Hochschulnetzwerk ohne Cloud-Einbindung betrieben werden soll, musste auf einen geringeren Ressourcenverbrauch geachtet werden.
Darüber hinaus benötigt Codetask aufgrund seiner überschaubaren Anzahl an Microservices aktuell noch keinen größeren Cluster.
Daher wurde sich gegen eine Standard-Kubernetes-Anwendung wie kubeadm oder Kubespray und für eine leichtgewichtige Kubernetes-Anwendung entschieden.

Dabei wurde die Entscheidung für K3s getroffen. Einerseits wurde ein k3s-Cluster parallel in einem anderen Projekt von Herrn Boger eingerichtet, andererseits bietet k3s auch darüber hinaus viele Vorteile.
Zum einen wurde k3s speziell entwickelt, um ein Kubernetes-Cluster in ressourcenbeschränkten Umgebungen auf Produktionsniveau auszuführen, beispielsweise auf IoT-Geräten.
Zum anderen hat k3s einen geringeren Einrichtungs- und Wartungsaufwand, da beim Installieren viele Komponenten, wie beispielsweise Traefik als Ingress-Controller, bereits integriert sind. Dadurch müssen weniger Komponenten separat installiert und konfiguriert werden.
Dadurch kann man sich mehr auf die eigentliche Anwendung konzentrieren statt auf die Clusteradministration.
Dennoch bleibt K3s einfach erweiterbar, unterstützt Hochverfügbarkeit, bietet eine gute Dokumentation und hat eine große Community.